\par\Needspace{6\baselineskip}\medskip
\hypertarget{native-097-article074}{}\pdfbookmark[2]{74.}{article-074}
\begin{center}74.\end{center}
Eine andere Aufgabe, \emph{aus dem geozentrischen Ort eines Himmelskörpers und der Lage seiner Bahnebene seinen heliozentrischen Ort auf der Bahn abzuleiten}, ist mit der vorhergehenden insofern verwandt, als auch sie vom Schnitt der zwischen der Erde und dem Himmelskörper gezogenen Geraden mit einer Ebene von gegebener Lage abhängt. Die Lösung erhält man am bequemsten aus den Formeln in Artikel 65. Dort hatten die Zeichen die folgende Bedeutung.

$L$ ist die Länge der Erde und $R$ ihr Abstand von der Sonne. Die Breite $B$ setzen wir $=0$ (denn der Fall, dass sie nicht $=0$ ist, lässt sich nach Artikel 72 leicht auf diesen Fall zurückführen), woraus $R'=R$ folgt. $l$ ist die geozentrische Länge des Himmelskörpers, $b$ seine Breite, $\Delta$ sein Abstand von der Erde, $r$ sein Abstand von der Sonne, $u$ das Breitenargument, $\AscendingNode$ die Länge des aufsteigenden Knotens und $i$ die Neigung der Bahn. Damit ergeben sich die Gleichungen
\[
\begin{aligned}
\mathrm{I.}\quad r\cos u-R\cos(L-\AscendingNode)&=\Delta\cos b\cos(l-\AscendingNode)\\[.5em]
\mathrm{II.}\quad r\cos i\sin u-R\sin(L-\AscendingNode)&=\Delta\cos b\sin(l-\AscendingNode)\\[.5em]
\mathrm{III.}\quad r\sin i\sin u&=\Delta\sin b.
\end{aligned}
\]
\noindent Multipliziert man Gleichung I mit $\sin(L-\AscendingNode)\sin b$, II mit $-\cos(L-\AscendingNode)\sin b$ und III mit $-\sin(L-l)\cos b$ und addiert die Produkte, so ergibt sich
\[
\begin{aligned}
&\cos u\sin(L-\AscendingNode)\sin b\\
&\quad-\sin u\cos i\cos(L-\AscendingNode)\sin b\\
&\quad-\sin u\sin i\sin(L-l)\cos b=0,
\end{aligned}
\]
\noindent woraus folgt
\[
\mathrm{IV.}\quad \operatorname{tang}u=\frac{\sin(L-\AscendingNode)\sin b}{\cos i\cos(L-\AscendingNode)\sin b+\sin i\sin(L-l)\cos b}\mathpunct{\cdot}
\]
\noindent \SourceLeaf{098}Multipliziert man nun I mit $\sin(l-\AscendingNode)$ und II mit $-\cos(l-\AscendingNode)$ und addiert die Produkte, so erhält man
\[
\mathrm{V.}\quad r=\frac{R\sin(L-l)}{\sin u\cos i\cos(l-\AscendingNode)-\cos u\sin(l-\AscendingNode)}\mathpunct{\cdot}
\]
\noindent Die Mehrdeutigkeit bei der Bestimmung von $u$ aus Gleichung IV löst sich durch Gleichung III von selbst auf. Sie zeigt, dass $u$ zwischen $0$ und $180^0$ oder zwischen $180^0$ und $360^0$ zu wählen ist, je nachdem, ob die Breite $b$ positiv oder negativ ist. Gilt hingegen $b=0$, so zeigt Gleichung V, dass $u=0$ oder $u=180^0$ zu setzen ist, je nachdem, ob $\sin(L-l)$ und $\sin(l-\AscendingNode)$ verschiedene oder gleiche Vorzeichen haben.

Die numerische Rechnung mit den Formeln IV und V lässt sich auf verschiedene Weise durch Einführung von Hilfswinkeln abkürzen. Zum Beispiel
\[
\begin{aligned}
\text{setzt man}\quad &\frac{\operatorname{tang}b\cos(L-\AscendingNode)}{\sin(L-l)}=\operatorname{tang}A,\\[.4em]
\text{so folgt}\quad &\operatorname{tang}u=\frac{\sin A\operatorname{tang}(L-\AscendingNode)}{\sin(A+i)}
\end{aligned}
\]
\[
\begin{aligned}
\text{setzt man}\quad &\frac{\operatorname{tang}i\sin(L-l)}{\cos(L-\AscendingNode)}=\operatorname{tang}B,\\[.4em]
\text{so folgt}\quad &\operatorname{tang}u=\frac{\cos B\sin b\operatorname{tang}(L-\AscendingNode)}{\sin(B+b)\cos i}\mathpunct{\cdot}
\end{aligned}
\]
\noindent Gleichung V nimmt ebenso eine übersichtlichere Form an, wenn man einen Winkel einführt, dessen Tangens $=\cos i\operatorname{tang}u$ oder $=\dfrac{\operatorname{tang}(l-\AscendingNode)}{\cos i}$ ist. So wie wir Formel V durch Kombination der Gleichungen I und II erhalten haben, gelangen wir durch Kombination der Gleichungen II und III zur folgenden Formel
\[
r=\frac{R\sin(L-\AscendingNode)}{\sin u(\cos i-\sin i\sin(l-\AscendingNode)\operatorname{cotang}b)}
\]
\noindent und ebenso durch Kombination der Gleichungen I und III zu dieser Formel
\[
r=\frac{R\cos(L-\AscendingNode)}{\cos u-\sin u\sin i\cos(l-\AscendingNode)\operatorname{cotang}b}\mathpunct{\cdot}
\]
\noindent Beide Formeln lassen sich ebenso wie V durch Einführung von Hilfswinkeln vereinfachen. Die aus dem Vorstehenden hervorgehenden Lösungen finden sich zusammengestellt und an einem Beispiel erläutert bei \textsc{von Zach}, \emph{Monatliche Correspondenz Bd. V. S. 540}\HistoricalNoteEight{Juni 1802. Band VI. S. 87.]}, weshalb wir hier auf eine weitere Ausführung verzichten. — Wird außer $u$ und $r$ auch der Abstand $\Delta$ gesucht, so lässt er sich aus Gleichung III bestimmen.
